\chapter{Landasan Teori}

%=========================================================%
%              TULIS LANDASAN TEORI DI SINI               %
% Jangan Lupa untuk Menghapus Contoh Tulisan di Bawah Ini %
%          (Termasuk Panduan di Setiap Subbabnya)         %
% Contoh Teks Dapat Di-comment dengan Tanda "%" di Depan  %
%   jika Ingin Mempertahankan Panduan di Dalam File Ini   %
%=========================================================%

% --- Tinjauan Pustaka atau Kajian Teori ---
\section{Kajian Teori}

Paparkan teori-teori, konsep, atau hukum dasar yang melandasi topik penelitian Anda. Bagian ini sebaiknya jangan hanya kumpulan definisi kamus, tapi harus berupa pembahasan mendalam yang menjadi rujukan utama untuk merencanakan tindakan perbaikan, pengembangan sistem, serta dasar analisis hasil pada bab berikutnya. Pastikan untuk memperdalam teori yang sebelumnya sudah disinggung sepintas di latar belakang, serta tambahkan konsep-konsep relevan lainnya yang mengikat variabel penelitian Anda. Setiap kutipan wajib mencantumkan sumber sitasi yang sah.



% --- Penelitian Terdahulu atau Kajian Komparasi Jurnal ---
\section{Penelitian Terdahulu yang Relevan}

Sajikan ulasan kritis terhadap penelitian-penelitian sejenis yang pernah dilakukan oleh peneliti lain sebelumnya (baik dari jurnal nasional maupun internasional). Jelaskan secara ringkas persamaan dan perbedaan antara penelitian mereka dengan penelitian yang sedang Anda lakukan (misalnya dari segi metode, objek, variabel, atau instrumen). Bagian ini ditujukan untuk menegaskan posisi penelitian Anda dan membuktikan bahwa karya Anda bebas dari plagiasi serta memiliki unsur kebaruan (\textit{novelty}). Sangat disarankan untuk merangkum bagian ini dalam bentuk tabel komparasi di akhir subbab. Tabel juga boleh saja dibuat dengan kertas \textit{landscape}.



% --- Kerangka Berpikir ---
\section{Kerangka Berpikir}

Uraikan alur logika pemikiran Anda dalam bentuk narasi analisis dan sintesis yang menggambarkan keterkaitan antara variabel atau komponen yang diteliti. Jelaskan bagaimana alur dari penemuan masalah di lapangan, proses pencarian data, penerapan metode, hingga rencana solusi yang diajukan untuk menemukan jawaban atas pertanyaan penelitian Anda. Selain dalam bentuk teks penjabaran, subbab ini wajib dilengkapi dengan bagan atau diagram alir (\textit{framework}) yang sistematis agar alur logika penelitian Anda mudah dipahami secara visual.



% --- Hipotesis ---
%
%     Dapat digunakan jika diperlukan dalam penelitian
%     kuantitatif
%
\section{Hipotesis Penelitian (Jika Diperlukan)}

Tuliskan kesimpulan atau jawaban sementara terhadap rumusan masalah penelitian yang kebenarannya akan diuji secara empiris melalui analisis data lapangan. Subbab ini umumnya wajib dicantumkan jika penelitian Anda menggunakan pendekatan kuantitatif atau eksperimen. Jika penelitian Anda bersifat deskriptif kualitatif atau murni pengembangan sistem (R\&D) tanpa pengujian hipotesis statistik, subbab ini dapat dihapus sepenuhnya dari dokumen, sekalian juga hapus kode \verb|\section{Hipotesis Penelitian}| atau di-\textit{comment} jika mau.